\uuid{cMQM}
\titre{Probabilité conditionnelle et indépendance}
\niveau{L1}
\module{Probabilité et statistique}
\chapitre{Loi, indépendance, loi conditionnelle}
\sousChapitre{Probabilité conditionnelle}
\theme{probabilité conditionnelle, indépendance, loi binomiale, urne, remise}
\auteur{Maxime Nguyen}
\datecreate{2026-07-11}
\organisation{AMSCC}
\difficulte{2}

\contenu{
	\texte{
		On dispose de deux urnes. La première contient 3 boules noires et 7 boules blanches. La seconde contient 5 boules noires et 5 boules blanches. 
		
		On choisit une urne au hasard et on tire successivement deux boules, avec remise, dans cette urne. On note les événements :
		\begin{itemize}
			\item $B1$ : \og on obtient une boule blanche au premier tirage\fg{}
			\item $B2$ : \og on obtient une boule blanche au second tirage\fg{}
		\end{itemize}
	}
	\question{Les événements $B1$ et $B2$ sont-ils indépendants ?}
	\indication{Conditionner par l'urne choisie. Les tirages sont indépendants sachant l'urne, mais il faut ensuite comparer $\prob(B_1\cap B_2)$ à $\prob(B_1)\prob(B_2)$.}
	\reponse{
		On note \(U_{1}\) (resp. \(U_{2}\)) l’événement « l’urne choisie est la première (resp. la seconde) ».  
		On a \(\Pr(U_{1})=\Pr(U_{2})=\dfrac12\).
		

		

		
		Dans l’urne 1 il y a \(7\) boules blanches sur \(10\) ; dans l’urne 2 il y a \(5\) boules blanches sur \(10\).  
		\[
		\Pr(B_{1}\mid U_{1})=\Pr(B_{2}\mid U_{1})=\frac{7}{10},\qquad 
		\Pr(B_{1}\mid U_{2})=\Pr(B_{2}\mid U_{2})=\frac{5}{10}=\frac12 .
		\]	
		
		
		Probabilité de \(B_{1}\) (et de \(B_{2}\)).
		\[
		\Pr(B_{1})=\Pr(U_{1})\Pr(B_{1}\mid U_{1})+\Pr(U_{2})\Pr(B_{1}\mid U_{2})
		=\frac12\cdot\frac{7}{10}+\frac12\cdot\frac12
		=\frac{7}{20}+\frac{5}{20}
		=\frac{12}{20}=\frac35 .
		\]
		Par symétrie, \(\Pr(B_{2})=\dfrac35\).
		
		
		
	
		\[
		\begin{aligned}
			\Pr(B_{1}\cap B_{2})
			&=\Pr(U_{1})\Pr(B_{1}\cap B_{2}\mid U_{1})+\Pr(U_{2})\Pr(B_{1}\cap B_{2}\mid U_{2})\\
			&=\frac12\left(\frac{7}{10}\right)^{2}+\frac12\left(\frac12\right)^{2}\\
			&=\frac12\cdot\frac{49}{100}+\frac12\cdot\frac{25}{100}
			=\frac{49+25}{200}
			=\frac{74}{200}
			=\frac{37}{100}.
		\end{aligned}
		\]
		
		

		Les événements \(B_{1}\) et \(B_{2}\) sont indépendants si et seulement si  
		\[
		\Pr(B_{1}\cap B_{2})=\Pr(B_{1})\Pr(B_{2}).
		\]
		
		Or  
		\[
		\Pr(B_{1})\Pr(B_{2})=\frac35\cdot\frac35=\frac{9}{25}=\frac{36}{100},
		\]
		tandis que  
		\[
		\Pr(B_{1}\cap B_{2})=\frac{37}{100}\neq\frac{36}{100}.
		\]
		
		Donc \(\Pr(B_{1}\cap B_{2})\neq\Pr(B_{1})\Pr(B_{2})\) et les événements \(B_{1}\) et \(B_{2}\) ne sont pas indépendants.  
		

	}
}
